% SOURCE docs/release/2026-10-29/source/docs/writing-repair/2026-10-29/src/8ed1637682836323.tex CHAR 0
\chapter{Longitudinal regimes, heterogeneous effects and continuous doses}
\label{v13:sw-regimes}

The detailed comparisons distinguish complete policy evaluation from components
that estimate heterogeneity or a dose curve. This chapter brings those choices
back to a small set of servicing and incentive requests. Its finite examples
show why the same data field or output name cannot determine which component
the bank needs. The preceding source studies supply the implementation limits;
the task here is to preserve them when composing the workflow.

\section{A sequence of decisions changes the software contract}

Servicing involves repeated decisions. A customer may receive a reminder, a
repayment plan and later a restructuring offer, with delinquency and contact
history changing between actions. An operation for a single binary treatment
does not become a longitudinal estimator by adding a time column. The target
must state which histories the policy observes, when actions occur, what ends
follow-up and how default or account closure changes the outcome.

Let $L_t$ be information available before action $A_t$, and let
$H_t=(\bar L_t,\bar A_{t-1})$ denote the observed history. A deterministic regime
$d$ supplies $A_t=d_t(H_t)$. For a terminal outcome $Y$, its value is
$V(d)=\mathbb E(Y^d)$. Under the required consistency, sequential
exchangeability and positivity conditions, a g-computation construction can
integrate an outcome regression against the covariate transition laws induced
by the intervention. Schematically,
\[
 V(d)=\int m(\bar l_T,\bar d_T)
             \prod_{t=0}^{T}p(l_t\mid\bar l_{t-1},\bar d_{t-1})
             \,d\bar l_T.
\]
The notation suppresses censoring and competing events, which must be specified
in an actual application. The product follows the intervention history; it is
not the factual joint distribution with treatment columns overwritten afterward.

The software comparison should therefore examine covariate transitions,
intervention callbacks, history updates, horizon conventions, simulation error
and sampling uncertainty. It should also compare parametric g-computation with
weighting and longitudinal targeted estimation on the same policy and population.
Disagreement over a default horizon or competing event can dominate any
difference in optimization speed.

\section{pygformula exposes the history and event choices}

The inspected \texttt{ParametricGformula} class accepts individual and time
identifiers, time-varying covariate names/types/models, custom histories, custom
fit/prediction functions and intervention descriptions. Its outcome options
include survival, continuous end-of-follow-up and binary end-of-follow-up.
Separate arguments address competing events, whether to treat them as censoring,
and a censoring model used for natural-course comparisons. \citep{v13_sw_pygformula}

For servicing, default and account closure should not be casually interchanged.
If default prevents later repayment, a default-free survival probability and a
cumulative incidence of repayment are different targets. If closing an account
merely hides future payment information, a censoring argument is needed. Asking
software to treat a competing event as censoring changes the scientific
experiment; the decision requires substantive justification and must survive
the library, CLI and agent interfaces.

The class documentation sets the number of bootstrap samples to zero by
default. The selected fit path invokes separate bootstrap work when samples
are requested. Thus the presence of bootstrap support does not imply that a
particular invocation produced sampling intervals. The helper review in
Section~\ref{sec:exec24-longitudinal-python} establishes whole-subject history
resampling and documents the failed-replicate handling discrepancy. That source
finding does not qualify arbitrary calendar dependence or the fitted intervals.
The adapter must retain the actual sampling unit and complete failure counts.
\citep{v13_sw_pygformula}

Simulation size and bootstrap size address different uncertainties. Conditional
on fitted transition laws, additional simulated trajectories reduce numerical
integration error. They do not provide new bank observations or reduce the
uncertainty in estimated transition parameters. A report should distinguish
the number of observed independent sampling units, simulated histories and
bootstrap replications. Treating simulated histories as additional observed
customers would greatly overstate precision.

The first useful adapter contract therefore preserves customer identity, event
and availability clocks, action timing, history definitions, horizon, competing
events and intervention semantics. It records whether uncertainty was requested
and actually obtained. Natural-course reproduction is a diagnostic: matching
factual outcomes does not prove that intervention predictions are valid.
Unsupported actions, missing positivity and impossible history combinations
should remain explicit, even if simulation can generate a numeric outcome.

The inspected MIT license supports considering reuse. Adoption still needs the
complete simulation and bootstrap paths, environment and dependency terms, and
matched comparisons with R implementations and other longitudinal estimators.
The retained source studies resolve selected simulation and bootstrap questions;
they do not replace those native qualification requirements.
Evaluating supplied regimes also does not establish that an optimal regime has
been learned. Policy selection and its evaluation are additional operations.

\section{grf: binary heterogeneity and continuous slopes are different targets}

The inspected \texttt{causal\_forest} documentation defines a binary-treatment
conditional effect and, for continuous treatment, describes the conditional
covariance ratio
\[
 \tau(x)=\frac{\operatorname{Cov}(Y,A\mid X=x)}
                    {\operatorname{Var}(A\mid X=x)}.
\]
The interface accepts nuisance predictions, clusters and equal cluster
weighting, honesty, subsampling, tuning, grouped trees for uncertainty and a
seed for its native random-number generator. Defaults include 2,000 trees,
honesty enabled and no tuning. These options define a specific statistical and
computational operation; they are not optional metadata for a reproducible
comparison. \citep{v13_sw_grf}

For binary treatment, an effect conditional on $X=x$ differs from a population
average unless it is integrated over the desired population. For continuous
treatment, a slope-like summary differs from a finite contrast
$\mu(a_1,x)-\mu(a_0,x)$. Calling both outputs ``uplift'' obscures the decision.
A bank deciding whether to increase an incentive from one specified amount to
another needs the latter contrast, or a justified approximation to it.

An exact project example makes the discrepancy visible. Suppose the stipulated
mean outcome, in CNY, is
\[
 \mu(a)=20a-3a^2,\qquad A\sim\operatorname{Unif}[0,2].
\]
This is a synthetic response curve, not a fitted bank relationship. Its
covariance slope is
\[
 \frac{\operatorname{Cov}(A,20A-3A^2)}{\operatorname{Var}(A)}
 =20-3\frac{\operatorname{Cov}(A,A^2)}{\operatorname{Var}(A)}=14.
\]
The finite change from dose one to dose two is instead
\[
 \mu(2)-\mu(1)=28-17=11.
\]
Exact moments of the uniform distribution verify both calculations. The
difference is a target difference, not a forest error. Multiplying the slope
by a unit change would overstate this stipulated finite benefit by three CNY.
Neither number includes incentive cost unless that cost is already part of the
outcome definition. \citep{v13_project_exact}

Cluster options require a similar question about the averaging measure. Giving
each branch equal total weight is different from giving each customer equal
weight when branch sizes differ. The chosen measure affects both the target
and the interpretation of heterogeneity. An adapter should require an explicit
population-weight policy and record the effective cluster structure; it should
not silently translate a customer population into an equal-branch population.

The R wrapper's seed, threading and native interface make an R/native adapter
a natural candidate. This inspection does not qualify the C++ training,
prediction or variance implementation. The observed repository GPL text also
requires a complete licensing and distribution review before a port or embedded
binary decision. A rewrite is not automatically preferable to a well-specified
boundary around the original implementation. \citep{v13_sw_grf}

\section{A generalized propensity score is a density component}

CausalGPS's inspected \texttt{estimate\_gps} function accepts continuous
exposure and covariates, a formula, normal or kernel density mode and a
SuperLearner library selection. In its normal branch, the fitted exposure mean
and a residual standard deviation enter a normal density. Its kernel branch
fits exposure and residual scaling and calls a separate density helper. Returned
fields include identity, exposure, generalized propensity score, fitted mean
and scale, and residuals. \citep{v13_sw_causalgps}

For continuous $A$, a generalized propensity score is a conditional density
$r(a,x)=f_{A\mid X}(a\mid x)$, not a probability in $[0,1]$. Its units depend
on the treatment scale. If $B=cA$ for $c>0$, then
\[
 f_{B\mid X}(b\mid x)=\frac1c f_{A\mid X}(b/c\mid x).
\]
Thus converting an incentive from CNY to hundreds of CNY changes the density
numerically. A bridge that normalizes every ``propensity'' into a binary
probability can corrupt continuous-treatment calculations. Exposure units and
the exact density evaluation points belong in the contract.

The observed-density function alone does not provide a dose-response estimator.
To compare feasible doses one needs densities or another identified construction
at the relevant counterfactual doses, an outcome model or weighting/matching
procedure, support diagnostics and an uncertainty argument.
Section~\ref{sec:exec24-dose-response-reuse} follows the density delegate,
transformation factors and downstream weighting. It identifies the selected
kernel branch's residual-density export and the conditions needed for an
original-scale density interpretation. Those restrictions must survive any
cross-package conversion. The retained source analysis establishes a component
and concrete qualification cases; it does not establish the complete causal
workflow.

A useful banking support diagnostic asks whether both proposed incentive levels
occur with adequate conditional support among the intended customers. The
normal model can assign positive density far beyond observed incentive choices;
mathematical positivity of a parametric density is not evidence of practical
overlap. Reporting density values together with observed range, local support
and model sensitivity is more informative than presenting a successful fit as
permission to extrapolate. The observed package declaration is GPL-3; complete
dependency and redistribution choices remain part of adoption review.

\section{causaldrf describes a full response surface, with different limits}

The inspected \texttt{hi\_est} documentation describes fitting a treatment
model, a GPS-adjusted outcome response surface, and predictions averaged over
customers at requested dose-grid values. It lists Normal, LogNormal, transformed
normal, Poisson, negative-binomial, Gamma and Binomial treatment models, and
returns fitted model and parameter records. This supports a documented
response-surface comparison; it is not an execution result or a review of every
internal estimator~\citep{v13_sw_causaldrf}.

Writing a fitted response surface as $\widehat m(a,r)$, the described averaging
operation is of the form
\[
 \widehat\mu(a)=\frac1n\sum_{i=1}^n
                      \widehat m\{a,\widehat r(a,X_i)\}.
\]
The GPS is evaluated at the proposed dose for each customer, rather than simply
reusing every customer's density at their observed dose. A finite effect then
requires a specified comparator, $\widehat\mu(a_1)-\widehat\mu(a_0)$.
A curve ordinate is an expected outcome, and labeling it as an effect without
the comparator misstates the answer.

The choice of dose grid also matters for decision support. If a business
constraint permits only three incentive levels, interpolation at arbitrary
levels has no immediate operating interpretation. If the optimum is selected
from a dense grid using the same fitted surface, uncertainty in that selection
is not the same as pointwise uncertainty at a fixed dose. A package comparison
should therefore distinguish fixed contrasts, whole curves and policy
optimization, and explain which one is actually supported.

The observed DESCRIPTION declares MIT plus a separate LICENSE file. The
initial documentation review did not include that accompanying notice or a
complete implementation review~\citep{v13_sw_causaldrf}. The later source
review recovered and inspected the notice, identifying the year and copyright
holder~\citep{v131_sw_causaldrf}. The appropriate decision remains a provisional
reference or adapter candidate: full implementation, interval behavior and
transitive dependency terms still require qualification. A recovered notice
is not a complete dependency inventory, and documented point estimation is
not evidence of qualified confidence intervals.

\section{A common servicing and incentive benchmark}

These components should be compared through a small set of declared banking
queries. One query fixes a two-period servicing regime and a terminal repayment
outcome. Another fixes a binary offer and asks for heterogeneity. A third fixes
two continuous incentive levels and asks for their net financial difference.
Each query has its own sampling unit, support requirements, feasible actions
and reference functional. A tool may be appropriate for one and unsupported
for another; that is useful information for architecture.

The comparison should retain a model that fits factual outcomes but fails an
intervention assumption, a weak-support population, a changed event definition
and a changed treatment unit. It should separately measure statistical error,
simulation error, numerical failure and computation. For longitudinal work,
the complete cost includes fitting transitions and histories, simulating
interventions, obtaining uncertainty and reporting diagnostics. Timing one
prediction call omits most of the workflow the bank would operate.

An eventual selection can combine components: symbolic identification, a
longitudinal engine, a heterogeneous-effect estimator and a dose-response
component can coexist behind one semantic contract. The common contract must
preserve their distinct assumptions and outputs. It should make unsupported
requests explicit rather than coercing every query into the nearest available
estimator.

